\documentclass{report}
\usepackage{amsmath,amssymb}
\setlength{\parindent}{0mm}
\setlength{\parskip}{1em}
\begin{document}
\begin{center}
\rule{6in}{1pt} \
{\large Matthias Bolten \\
{\bf Local {Fourier} Analysis for Block Smoothers in Multigrid Methods}}

Bergische Universit\"at Wuppertal \\ Fachbereich Mathematik und Naturwissenschaften \\ 42097 Wuppertal \\ Germany
\\
{\tt bolten@math.uni-wuppertal.de}\\
Karsten Kahl\end{center}


On parallel computers and nowadays computer architectures block smoothers
possess the advantage of a high computation/communication and/or a
computation/memory access ratio, especially when compared to traditional
point smoothers. Therefore, multigrid methods should benefit from the use
of block smoothers regarding time to solution.

While block smoothers like block Jacobi are used in current parallel
implementations of multigrid methods, they have not been studied
extensively. We extended the traditional Local {Fourier} Analysis
described in the literature in detail e.g.~in the book of Wienands and
Joppich for the case of block smoothers. Unlike the previous study of
overlapping smoothers in the recent paper by MacLachlan and Oosterlee
(Num. Lin. Alg. Appl., 18) our approach is based on a decomposition of
the infinite Toeplitz matrices and considers matrix-valued generating
symbols instead of scalar generating symbols.

We present our approach and provide numerical examples for different
block smoothers, including smoothers with and without overlap and
different coarsening ratios and block sizes. Additionally, we consider
different solution strategies for the block solves, including direct
solvers as well as iterative methods.


\end{document}
